Source : calculs des auteurs, \texttt{note8p/calculs/is\_reforme.py} (convention de Devereux et Griffith).\par Équipement neuf d'une durée d'usage de 7~ans (plan actuel pour moitié linéaire, pour moitié dégressif), financé par fonds propres ; rendement réel exigé 5~\%, inflation 2~\%, dépréciation 14~\% par an, sans impôt personnel. Déduction de 50~\% : la moitié du prix déduite l'année de l'achat, le reste amorti normalement. Coûts bruts, hors fin du taux réduit : 75~\% de l'investissement éligible fait par des sociétés imposables ; coût permanent une fois la transition passée, en euros de 2027 (il croît ensuite avec l'investissement). Financé par emprunt, le même équipement a aujourd'hui un taux de $-24$~\% ; $-33$~\% avec la combinaison.
